%! TEX program = xelatex

\documentclass[withoutpreface,bwprint]{cumcmthesis}

\usepackage[most]{tcolorbox}
\usepackage{tikz}
\usepackage{cumcm2026}
\usepackage[numbers,sort&compress]{natbib}
\usepackage{url}

\title{LATEX_TEMPLATE_TITLE}

\begin{document}

\maketitle

% ============================================================
% 摘要
% ============================================================

\begin{abstract}

  LATEX_TEMPLATE_ABSTRACT

  \keywords{LATEX_TEMPLATE_KEYWORDS}

\end{abstract}

% ============================================================
% 正文
% ============================================================

\section{问题重述}

LATEX_TEMPLATE_BODY

% 公式规范（硬性）：
%   - 行间公式必须用编号环境 equation / align（公式居中、编号右对齐，LaTeX 自动处理）；
%   - 禁止 \[...\]、equation*、align*、手写 \tag；
%   - 正文用 \eqref{eq:...} 引用；公式出现后解释每个符号。

\section{问题分析}

LATEX_TEMPLATE_BODY

\section{模型假设与符号说明}

LATEX_TEMPLATE_BODY

\section{数据预处理}

LATEX_TEMPLATE_BODY

% 逐问建模与求解：按题目子问题数量展开，每问一节；子问题只有两个时删除多余小节。
% 每问包含：建模思路 → 数学表达 → 求解方法 → 结果 → 图表 + 紧随其后的分析。

\section{问题一建模与求解}

LATEX_TEMPLATE_BODY

\section{问题二建模与求解}

LATEX_TEMPLATE_BODY

\section{问题三建模与求解}

LATEX_TEMPLATE_BODY

% 如果题目没有第四问，则删除该节
\section{问题四建模与求解}
LATEX_TEMPLATE_BODY

\section{模型的分析与检验}

LATEX_TEMPLATE_BODY

\section{模型的评价与推广}

LATEX_TEMPLATE_BODY

\CumcmAIUsed{[简要用途, 如语言润色, 代码调试等]}

\bibliographystyle{gbt7714-numerical}
\bibliography{refs}

% 附录：支撑材料文件列表

\begin{appendices}

  \section{附录}

  LATEX_TEMPLATE_BODY

\end{appendices}

\end{document}